
\subsubsection{5.1 Infrastructure Validation}

\#\#\#\# 5.1.1 Tier Viability (RQ1)

K=5 solution generation achieves perfect coverage (1.0000) for all 542 problems across all three models. Tier stratification results:

\begin{table*}[htbp]
\centering
\resizebox{\dimexpr 2\columnwidth+\columnsep\relax}{!}{%
\begin{tabular}{lllll}
\toprule
Model & Benchmark & n\_hard & n\_easy & Viable \\
\midrule
Llama3-8B & HumanEval+ & 78 & 39 & Yes \\
Llama3-8B & MBPP+ & 150 & 128 & Yes \\
CodeLlama-7B & HumanEval+ & 142 & 0 & No (n\_easy=0) \\
CodeLlama-7B & MBPP+ & 199 & 37 & Yes \\
DeepSeek-6.7B & HumanEval+ & 68 & 24 & Yes \\
DeepSeek-6.7B & MBPP+ & 105 & 176 & Yes \\
\bottomrule
\end{tabular}
}
\end{table*}

CodeLlama-7B achieves n\_easy=0 on HumanEval+, reflecting that this base model rarely achieves pass@1 >= 0.6 on HumanEval+ problems without instruction tuning. MBPP+ serves as the primary benchmark for CodeLlama's easy tier analysis. The gate passes with 5/6 pairs viable.

\#\#\#\# 5.1.2 P(True) Non-Degeneracy (RQ3)

P(True) extraction produces non-degenerate confidence distributions for all three models across 5,730 (problem, solution) pairs:

\begin{table*}[htbp]
\centering
\resizebox{\dimexpr 2\columnwidth+\columnsep\relax}{!}{%
\begin{tabular}{lllll}
\toprule
Model & mean(c) & std(c) & r(c, correctness) & p-value \\
\midrule
Llama3-8B & 0.4989 & 0.0669 & +0.197 & < $10^-18$ \\
CodeLlama-7B & 0.3682 & 0.0618 & +0.144 & < $10^-10$ \\
DeepSeek-6.7B & 0.6480 & 0.0781 & -0.046 & 0.048 \\
\bottomrule
\end{tabular}
}
\end{table*}

All models satisfy std(c) > 0.05. The confidence-correctness correlation is statistically significant but weak (r = 0.14-0.20 for Llama3 and CodeLlama). DeepSeek shows a near-zero negative correlation (-0.046), indicating that its higher mean confidence does not align positively with correctness at the instance level.

\#\#\#\# 5.1.3 Tier Consistency (RQ2)

Cross-architecture Jaccard similarity of hard-tier assignments:

\begin{table}[htbp]
\centering
\resizebox{\columnwidth}{!}{%
\begin{tabular}{lll}
\toprule
Model Pair & Jaccard & Gate (> 0.30) \\
\midrule
Llama3 and CodeLlama & 0.546 & Pass \\
Llama3 and DeepSeek & 0.531 & Pass \\
CodeLlama and DeepSeek & 0.456 & Pass \\
\bottomrule
\end{tabular}
}
\end{table}

All pairs exceed the 0.30 threshold. 133 of 542 problems (24.5\%) are hard for all three models, forming an architecture-independent difficulty core.

\subsubsection{5.2 Main Result: Architecture-Dependent DELTA\_ECE (RQ4)}

\begin{table*}[htbp]
\centering
\resizebox{\dimexpr 2\columnwidth+\columnsep\relax}{!}{%
\begin{tabular}{lllllllll}
\toprule
Model & Architecture & n\_hard & n\_easy & ECE(hard) & ECE(easy) & DELTA\_ECE & 95\% CI & p-value \\
\midrule
Llama3-8B & General & 228 & 167 & 0.4887 & 0.4852 & +0.0034 & [-0.0074, +0.0133] & 0.256 \\
CodeLlama-7B & Code-adapted & 341 & 37 & 0.3659 & 0.6149 & -0.2490 & [-0.2589, -0.2391] & 1.000 \\
DeepSeek-6.7B & Code-specialized & 173 & 200 & 0.6565 & 0.3586 & +0.2979 & [+0.2849, +0.3115] & 0.000 \\
\bottomrule
\end{tabular}
}
\end{table*}

\textbf{Pre-registered gate result: 1/3 models pass (requires >= 2/3). GATE FAIL.}

Three qualitatively distinct patterns emerge:

\textbf{DeepSeek-Coder (DELTA\_ECE = +0.298).} The code-specialized model shows the expected pattern: hard problems are substantially more miscalibrated than easy ones. The bootstrap CI [0.285, 0.312] is entirely positive.

\textbf{Llama3-8B (DELTA\_ECE approximately 0).} The general-purpose model shows no meaningful stratification effect. ECE(hard) approximately equals ECE(easy) approximately equals 0.49, and the 95\% CI includes zero. Both tiers are substantially miscalibrated, but the miscalibration is uniform across difficulty levels.

\textbf{CodeLlama-7B (DELTA\_ECE = -0.249).} The code-adapted model shows the opposite of the expected pattern. ECE(easy) = 0.615 exceeds ECE(hard) = 0.366, meaning easy problems are more miscalibrated. The CI [-0.259, -0.239] is entirely negative.

\begin{figure}[htbp]
\centering
\includegraphics[width=\columnwidth]{figures/fig1_delta_ece_gate.png}
\caption{DELTA\_ECE per model with bootstrap CI}
\label{fig:fig1_delta_ece_gate}
\end{figure}

\begin{figure}[htbp]
\centering
\includegraphics[width=\columnwidth]{figures/fig2_reliability_diagrams.png}
\caption{Reliability diagrams by tier per model}
\label{fig:fig2_reliability_diagrams}
\end{figure}

\subsubsection{5.3 Temperature Scaling Analysis}

\begin{table*}[htbp]
\centering
\resizebox{\dimexpr 2\columnwidth+\columnsep\relax}{!}{%
\begin{tabular}{lllll}
\toprule
Model & T* & DELTA\_ECE (pre) & DELTA\_ECE (post) & Direction preserved \\
\midrule
Llama3-8B & 1.163 & +0.0034 & -0.1371 & No (inverted) \\
CodeLlama-7B & 3.951 & -0.2490 & -0.8099 & Yes (worsened) \\
DeepSeek-6.7B & 1.210 & +0.2979 & +0.0728 & Yes (reduced but persists) \\
\bottomrule
\end{tabular}
}
\end{table*}

Global temperature scaling does not correct architecture-dependent DELTA\_ECE direction. CodeLlama's T* = 3.95 is approximately 3x the correction needed for the other models, indicating substantial global confidence inflation from code fine-tuning.

\begin{figure}[htbp]
\centering
\includegraphics[width=\columnwidth]{figures/fig3_temperature_scaling.png}
\caption{Temperature scaling effect}
\label{fig:fig3_temperature_scaling}
\end{figure}

\subsubsection{5.4 M-Sensitivity}

DELTA\_ECE values are stable across M in {10, 15, 20} for all models (Llama3: +0.00344; CodeLlama: -0.24899; DeepSeek: +0.29789 -- unchanged across bin counts), ruling out bin-count artifacts.

\subsubsection{5.5 Cross-Architecture Difficulty Overlap}

\begin{figure}[htbp]
\centering
\includegraphics[width=\columnwidth]{figures/jaccard_similarity_bars.png}
\caption{Jaccard similarity across model pairs}
\label{fig:jaccard_similarity_bars}
\end{figure}

\begin{figure}[htbp]
\centering
\includegraphics[width=\columnwidth]{figures/consensus_hard_pie.png}
\caption{Consensus hard problem distribution}
\label{fig:consensus_hard_pie}
\end{figure}

The Jaccard similarity of hard-tier assignments ranges from 0.456 to 0.546 across model pairs. 133 problems (24.5\% of EvalPlus) are universally hard across all three architectures.
